%% kapa: exam=1 round=31 year=2020 period=2 form=A subject=appraisal_law subject_ko=감정평가관계법규 seq=35 number=35
%% source: Q-Net 감정평가사 기출문제 (cst00309) artlSeq=5207219 file=extracted/1008/5207219/A형/A형/2020년도 제31회 감정평가사 1차 시험_2교시_A형.pdf
%% conversion: kordoc 4.24.51 -> tools/gen1.py (auto); review=proofread-sonnet; flags=none
\begin{question}{35}
부동산등기법령상 ‘권리에 관한 등기’에 관한 설명으로 옳은 것은?
\begin{choices}
  \item 권리자가 2인 이상인 경우에는 권리자별 지분을 기록하여야 하고 등기할 권리가 총유(總有)인 때에는 그 뜻을 기록하여야 한다.
  \item 등기원인에 권리의 소멸에 관한 약정이 있을 경우 신청인은 그 약정에 관한 등기를 신청할 수 있다.
  \item 등기관이 소유권 외의 권리에 대한 처분제한 등기를 할 때 등기상 이해관계 있는 제3자의 승낙이 없으면 부기로 할 수 없다.
  \item 등기관이 환매특약의 등기를 할 때 매매비용은 기록하지 아니한다.
  \item 등기관이 소유권보존등기를 할 때 등기원인과 그 연월일을 기록하여야 한다.
\end{choices}
\end{question}
\answer{2}
